
We demonstrate that Direct Preference Optimization and attribute conditioning can be jointly optimized in a single training run through gradient-compatible multi-task learning. Gradient angle measurements revealed a mean of 78.5° (SD: 12.8°) between task gradients, well below the 120° catastrophic interference threshold, with zero measurements exceeding this threshold across 100 proof-of-concept training steps. Both DPO loss (5.8\% reduction) and attribute loss (21.3\% reduction) decreased monotonically, confirming convergence without destructive task conflict.

The jointly trained model achieved 54\% preference win rate (maintaining 94\% of standalone DPO baseline performance) while simultaneously achieving 65\% attribute steering accuracy (exceeding random chance baselines by 45 percentage points), both surpassing proof-of-concept feasibility thresholds. Linear probing revealed 100\% accuracy in classifying preferences from hidden states, demonstrating that joint training creates shared representations encoding task-relevant information.

We establish gradient compatibility as a quantitative design principle for multi-objective LLM alignment. While our proof-of-concept experiments validate feasibility at 100-step scale, full-scale deployment and extension to N>2 objectives remain open questions. Nevertheless, the gradient angle measurement methodology---testing whether angle($\nabla Obj1, \nabla Obj2)$ < 120°---provides a transferable criterion for predicting joint optimization feasibility before expensive training runs.

Future work includes full-scale validation at 15,000 training steps with loss weight ablation, training sequential baselines to test emergent benefits, integrating real attribute labels to measure representation disentanglement, and extending to N=3 objectives (Constitutional AI + DPO + Attributes) to test whether compatibility scales to multi-way combinations.
